\section{Bản chất và các phương pháp sản xuất giá trị thặng dư}

\subsection{Bản chất giá trị thặng dư}

\subsubsection{Góc nhìn của C. Mác}

Như đã đề cập ở trên, giá trị thặng dư tồn tại xuất phát từ chính hao phí sức lao động sống của người công nhân, tức là không do máy móc sinh ra. Khi giả định xã hội chỉ có hai giai cấp gồm tư sản và công nhân thì giá trị thặng dư trong nền kinh tế thị trường tư bản chủ nghĩa mang bản chất kinh tế - xã hội \cite[trang 89]{gt}, phản ánh quan hệ giai cấp và sự bóc lột. Nhà tư bản làm giàu dựa trên việc thuê mướn sức lao động của công nhân.\\\\
Trong khi các bậc tiền bối như Adam Smith hay David Ricardo chỉ mới dừng lại ở việc quan sát "lợi nhuận" như một hiện tượng bề mặt của thị trường, C. Mác đã thực hiện một bước nhảy vọt khi đi sâu vào bản chất của sản xuất để phát hiện ra bản chất của giá trị thặng dư. Sự phát hiện khoa học của C. Mác chỉ ra rằng nhà tư bản bóc lột công nhân mà không hề vi phạm quy luật trao đổi ngang giá. Nghĩa là, nhà tư bản mua hàng là hóa sức lao động đúng giá trị của nó, nhưng giá trị thặng dư vẫn được sinh ra trong quá trình sử dụng lao động đó. Vì lẽ đó, mục đích nghiên cứu của Mác không chỉ dừng lại ở việc giải thích cách thức vận hành của các thực thể kinh tế, mà quan trọng hơn, là lột trần quan hệ giai cấp bóc lột được che đậy tinh vi, từ đó khẳng định tính tất yếu của sự thay thế các quan hệ sản xuất lỗi thời. Ngày nay, quan hệ bóc lột này vẫn tồn tại nhưng dưới các hình thức tinh vi và văn minh hơn so với thế kỷ XIX.

\subsubsection{Thước đo bản chất và quy mô bóc lột}

Trong nền kinh tế trư bản chủ nghĩa, các nhà tư bản luôn nỗ lực để tối đa hóa giá trị thặng dư. Để đo lường cường độ và phạm vi của sự chiếm đoạt, lý luận Mác-xít thiết lập hai đại lượng định lượng khách quan đó là tỷ suất và khối lượng giá trị thặng dư.\\

\textbf{Tý suất giá trị thặng dư}\\\\
Đại lượng được định nghĩa như sau: “Tỷ suất giá trị thặng dư là tỷ lệ phần trăm giữa giá trị thặng dư và tư bản khả biến để sản xuất ra giá trị thặng dư đó.” \cite[trang 90]{gt} Công thức tính là:

$$m'=\frac{m}{v}\times100\%$$

Trong đó:
\begin{itemize}
    \item $m’$ là tỷ suất giá trị thăng dư;
    \item $m$ là giá trị thặng dư;
    \item $v$ là tư bản khả biến.
\end{itemize}
Ngoài ra, tỷ suất giá trị thặng dư còn được tính theo công thức sau:
$$m'=\frac{t'}{t}\times100\%$$
Trong đó:
\begin{itemize}
    \item $t’$ là thời gian lao động thặng dư - phần thời gian làm việc không công cho nhà tư bản, tạo ra giá trị thặng dư $(m)$;
    \item $t$ là thời gian lao động tất yếu - phần thời gian người công nhân tạo ra giá trị bù đắp lại giá trị sức lao động $(v)$.
\end{itemize}
Dựa vào công thức trên, chỉ số này phản ánh cường độ bóc lột, tức là nhà tư bản đã "vắt kiệt" bao nhiêu phần trăm từ một đơn vị sức lao động. Nó cho thấy mức độ hiệu quả của việc chiếm đoạt lao động không công trên mỗi công nhân. Như vậy, tỷ suất giá trị thặng dư là biểu hiện chính xác của trình độ tư bản bóc lột sức lao động, nhà tư bản bóc lột người công nhân.

\textbf{Khối lượng giá trị thặng dư}\\\\
Đại lượng được định nghĩa như sau: “Khối lượng giá trị thặng dư là lượng giá trị thặng dư bằng tiền mà nhà tư bản thu được” 
\cite[trang 91]{gt}, công thức tính là:
$$ M = m'\times V$$

Trong đó:
\begin{itemize}
    \item $M$ là khối lượng giá trị thặng dư;
    \item $m'$ là tỷ suất giá trị thặng dư;
    \item $V$ là tổng tư bản khả biến.
\end{itemize}
	Đại lượng này phản ánh quy mô bóc lột, cho thấy tổng giá trị mà nhà tư bản thu được từ toàn bộ đội ngũ lao động làm thuê. Một doanh nghiệp có thể có tỷ suất giá trị thặng dư thấp hơn nhưng lượng giá trị thặng dư khổng lồ nếu họ thuê mướn một số lượng nhân công cực lớn.

\subsection{Các phương pháp sản xuất giá trị thặng dư}

\subsubsection{Sản xuất giá trị thặng dư tuyệt đối}

“Giá trị thặng dư tuyệt đối là giá trị thặng dư thu được do kéo dài ngày lao động vượt quá thời gian lao động tất yếu, trong khi năng suất lao động, giá trị sức lao động và thời gian lao động tất yếu không thay đổi” \cite[trang 91]{gt}. Để dễ hình dung ta lấy ví dụ trang 91 trong Giáo trình Kinh tế chính trị Mác Lê-nin:\\\\
Nếu xét một ngày lao động 8 giờ, trong đó 4 giờ là tất yếu $(v)$ và 4 giờ là thặng dư $(m)$, khi đó $m' = 100\%$. Nếu nhà tư bản kéo dài ngày làm việc lên 10 giờ mà không tăng lương (giữ $v$ là 4 giờ), thời gian lao động thặng dư tăng lên 6 giờ, lúc này $m'$ tăng vọt thành $150\%$. Ví dụ này minh chứng rằng trình độ bóc lột luôn có xu hướng gia tăng không ngừng thông qua việc kéo dài ngày lao động hoặc tăng năng suất lao động để giảm thời gian tất yếu. Nhà tư bản thèm khát vô hạn đối với lao động thặng dư nên tìm mọi cách kéo dài nó không bị pháp luật hạn chế. Kéo dài ngày lao động quá mức và sử dụng lao động trẻ em, vắt kiệt tinh thần và thể xác người lao động. Marx gọi cách khai thác như vậy là bóc lột sức lao động.

\subsubsection{Sản xuất giá trị thặng dư tương đối}

“Giá trị thặng dư tương đối là giá trị thặng dư thu được nhờ rút ngắn thời gian lao động tất yếu; do đó kéo dài thời gian lao động thặng dư trong khi độ dài ngày lao động không thay đổi hoặc thậm chí rút ngắn” \cite[trang 92]{gt}. Để dễ hình dung ta lấy ví dụ trang 93 trong Giáo trình Kinh tế chính trị Mác Lê-nin:\\\\
Ngày lao động 8 giờ, với 4 giờ lao động tất yếu, 4 giờ lao động thặng dư, tỷ suất giá trị thặng dư là $100\%$. Nếu giá trị sức lao động giảm khiến thời gian lao động tất yếu rút xuống còn 2 giờ thì thời gian lao động thặng dư sẽ là 6 giờ. Khi đó, $m’ = 300\%$. Nếu ngày lao động giảm xuống còn 6 giờ nhưng giá trị sức lao động giảm khiến thời gian lao động tất yếu rút xuống còn 1 giờ thì thời gian lao động thặng dư sẽ là 5 giờ. Khi đó: $m’ = 500\%$.\\\\
Để nhà tư bản thu được nhiều giá trị thặng dư hơn mà không cần kéo dài ngày lao động, họ phải tìm cách hạ thấp giá trị sức lao động. Cách thực hiện ở đây chính là tăng năng suất lao động xã hội trong các ngành sản xuất tư liệu sinh hoạt (nhu yếu phẩm) và các ngành chế tạo máy móc sản xuất ra các tư liệu sinh hoạt đó. Việc này giúp hàng hoá do các xí nghiệp ấy sản xuất ra có giá trị cá biệt thấp hơn giá trị xã hội, thu được một số giá trị thặng dư trội hơn so với các xí nghiệp khác (giá trị thặng dư siêu ngạch). Máy móc biến một bộ phận tư bản trước kia là khả biến thành máy móc, tức thành tư bản bất biến. Vì vậy, sử dụng máy móc vào sản xuất là phương pháp hữu hiệu nhất để rút ngắn thời gian lao động cần thiết, kéo dài thời gian lao động thặng dư.